% =====================================================================
%  File: sections/3_5_thiet_ke_benchmark.tex
%  Mục đích: Phần 3.5 - Thiết kế benchmark
% =====================================================================

\subsection{Thiết kế benchmark}
\label{sec:thiet-ke-benchmark}

Benchmark của Raptix được thiết kế để phục vụ hai mục đích song song: (i) so sánh Raptix với các hệ thống red team tự động hiện có, và (ii) làm khung non-regression cho quá trình phát triển, đảm bảo mỗi thay đổi kiến trúc thực sự mang lại cải tiến chứ không chỉ dịch chuyển chi phí từ trục này sang trục khác. Đây là điểm khác biệt căn bản so với cách đánh giá truyền thống chỉ dựa trên solve rate.

\paragraph{Động cơ thiết kế}

Nếu chỉ đo ``hệ thống tìm được bao nhiêu lỗ hổng'' hoặc ``chiếm được bao nhiêu host'', một agent có thể đạt điểm cao bằng cách gọi rất nhiều tool và tiêu thụ một lượng lớn token, mà không thực sự thông minh hơn. Ngược lại, nếu chỉ đo chi phí, hệ thống có thể tối ưu bằng cách không làm gì cả. Vì vậy, mỗi run benchmark cần thu đồng thời cả \emph{outcome} (kết quả đạt được) và \emph{cost} (chi phí bỏ ra), đồng thời chạy lặp nhiều lần để đo độ ổn định. Ba trục Effectiveness, Efficiency và Reliability chính là cơ sở để cân bằng hai khía cạnh này.

\subsubsection*{Ba trục đánh giá chính}

\begin{table}[H]
\centering
\caption{Ba trục đánh giá chính và chỉ số đo lường}
\label{tab:ba-truc-danh-gia}
\small
\renewcommand{\arraystretch}{1.3}
\begin{tabular}{|p{2.8cm}|p{6cm}|p{6.5cm}|}
\hline
\textbf{Trục} & \textbf{Câu hỏi trung tâm} & \textbf{Chỉ số đo lường} \\
\hline
Effectiveness & Hệ thống có đạt được mục tiêu không? & Task success rate, số verified findings, số flag/objective đạt được, mức độ hoàn thành attack-chain. \\
\hline
Efficiency & Hệ thống đạt mục tiêu với chi phí bao nhiêu? & Token tiêu thụ, số model calls, số tool calls, wall-clock time, số vòng retry, coordination overhead. \\
\hline
Reliability & Kết quả có lặp lại được không? & Tỷ lệ thành công qua nhiều run, variance, khả năng recovery sau timeout/failure, state consistency. \\
\hline
\end{tabular}
\end{table}

Ba trục này có quan hệ đánh đổi với nhau. Một thay đổi kiến trúc có thể tăng Effectiveness nhưng giảm Efficiency (ví dụ: thêm nhiều agent con để khám phá song song), hoặc tăng Reliability nhưng giảm Efficiency (ví dụ: thêm verifier độc lập cho mọi finding). Vì vậy, một thay đổi chỉ được coi là cải tiến toàn diện khi \emph{không} làm suy giảm nghiêm trọng hai trục còn lại.

\subsubsection*{Các bộ benchmark sử dụng}

Khung đánh giá của Raptix lựa chọn năm benchmark và môi trường đánh giá nhằm bao phủ các nhóm capability khác nhau, từ web exploitation tới Active Directory và long-horizon operations:

\begin{table}[H]
\centering
\caption{Các benchmark và môi trường đánh giá, cùng capability mục tiêu}
\label{tab:benchmark-suite}
\small
\renewcommand{\arraystretch}{1.3}
\begin{tabular}{|p{3.4cm}|p{5.5cm}|p{6.5cm}|}
\hline
\textbf{Benchmark} & \textbf{Capability chính} & \textbf{Vai trò trong suite} \\
\hline
WebExploitBench~\cite{agentcyberrange} & Web exploitation trên target cô lập & Đo khả năng tìm và xác minh lỗ hổng web thực tế (phù hợp pha Initial Foothold). \\
\hline
PostExploitBench~\cite{agentcyberrange} & Multi-host, post-exploitation, pivot & Đo mission nhiều bước và phụ thuộc trạng thái (phù hợp pha Network Propagation). \\
\hline
GOAD~\cite{goad} & Windows / Active Directory & Đo enumeration, privilege path và lateral movement trong AD lab. \\
\hline
Nilgiri~\cite{nilgiri} & Long-horizon cyber range & Đo khả năng duy trì kế hoạch và state qua chuỗi nhiệm vụ dài. \\
\hline
InfraGauntlet~\cite{infragauntlet} & Linux/network, non-HTTP services & Bổ sung pure infrastructure, false-positive control và pivoting. \\
\hline
\end{tabular}
\end{table}

Việc chọn nhiều bộ benchmark khác nhau nhằm tránh thiên lệch về một nhóm lỗ hổng dễ đo như web injection. Một hệ thống có thể đạt điểm cao trên WebExploitBench nhưng thất bại hoàn toàn trên GOAD hoặc Nilgiri nếu kiến trúc của nó không được thiết kế cho mission dài hoặc AD attack. Do đó, việc đánh giá Raptix cần được thực hiện trên \emph{toàn bộ} suite, không chỉ trên một bộ đơn lẻ.

\subsubsection*{Chỉ số bổ trợ cho mô hình hai pha}

Vì Raptix dùng mô hình hai pha, benchmark ghi thêm: (i) \textbf{thời điểm chuyển pha} (số bước/tool call trước khi chuyển từ Initial Foothold sang Network Propagation); (ii) \textbf{số attack path đã thử trong AKG}; (iii) \textbf{tỷ lệ credential reuse thành công} (đo hiệu quả truy vấn \texttt{CAN\_AUTH\_AS} và \texttt{PIVOT\_VIA}); (iv) \textbf{số nhánh đã loại trừ} (đo cơ chế tránh lặp hành động).

\subsubsection*{Pin baseline và quy trình non-regression}

Trước mỗi vòng tối ưu, cần pin: commit hash, model/provider/temperature, toolset (tên + version), cấu hình benchmark và giới hạn tài nguyên. Sau đó chạy baseline ít nhất 5 lần, áp dụng thay đổi, chạy lại cùng điều kiện. Do quy mô pilot nhỏ (5 run), báo cáo chủ yếu sử dụng thống kê mô tả: tỷ lệ thành công với Wilson interval cho chỉ số nhị phân; median và khoảng dao động cho chỉ số liên tục. Kiểm định suy diễn chỉ được áp dụng khi số lần chạy đủ lớn theo thiết kế thí nghiệm và phân tích power, đồng thời lựa chọn phép kiểm định phù hợp với loại metric. Thay đổi chỉ được giữ khi cải thiện ít nhất một trục mà không làm suy giảm nghiêm trọng hai trục còn lại.

